
\شروع‌مسئله [15]
در کشور شایان تعدادی شهر وجود دارد که با تعدادی جاده دوطرفه با هم در ارتباط هستند و میدانیم بین هر دو شهری دقیقا یک مسیر وجود دارد.\\*
شایان می خواهد هر دو شهری که  به او داده می شود باید فاصله میان آن دو شهر را بگوید اما نمی تواند.
به او کمک کنید تا فاصله دو شهر را بگوید.

\متن‌سیاه{ورودی:}\\
در خط اول به ترتیب دو عدد  $n,q$ داده می شود که بترتیب تعداد شهر ها و تعداد پرسش فاصله است.\\*
سپس در $n-1$ خط سه عدد $i,j,w$ داده می شود که نشان می دهد بین دو شهر با شماره $i,j$ مسیری به طول $w$ وجود دارد.\\*
سپس در $q$ خط دو عدد $i,j$ داده می شود که پرسش فاصله بین دو شهر با آن شماره است.\\*
(تضمین می شود گراف داده شده شرایط اولیه مسئله را دارد)

\متن‌سیاه{خروجی:}\\
جواب $q$ پرسش را در $q$ خط زیر هم چاپ کنید.

\متن‌سیاه{محدودیت‌ها:}\\
$1 \leq n \leq 4 * 10^4$ \\
$1 \leq q \leq 10^6$

ورودی و خروجی نمونه:

\begin{latin}
	\begin{tabularx}{\textwidth}{|X | X|}
		\hline
		\textbf{stdin} &
		\textbf{stdout}\\
		\hline
		6 4
		
		1 2 5
		
		1 3 4
		
		2 4 3
		
		2 5 2
		
		3 6 1
		
		1 6
		
		4 5
		
		2 6
		
		4 6&
		
		5 
		
		5 
		
		10 
		
		13\\
		\hline
		
	\end{tabularx}
\end{latin}


\solution{
یک راس را در نظر می گیریم و درخت را از آن راس ریشه دار می کنیم. سپس فاصله هر راسی تا آن راس ریشه را هم حساب می کنیم. این دو کار در $O(n)$ قابل انجام است.

سپس یک آرایه بصورت بازگشتی میسازیم که پیدا کردن نزدیک ترین جد مشترک دو راس بشود همان مسئله پیدا کردن مینیمم در آرایه بین اندیس $i$ ام تا $j$ام. این کار هم در  $O(n)$قابل انجام است.

سپس چون آپدیت نداریم از روش توان های دو در $RMQ$ استفاده می کنیم و یک پیش پردازش در $O(n\log n)$ انجام می دهیم و برای جواب هر پرسش هم در $O(1)$ جواب می دهیم و حافظه الگوریتم هم $O(n\log n)$ است.

پس کل الگوریتم دز $O(n\log n + q)$ قابل انجام است.

}

